
\section{An upper bound of $\widetilde{O}(n^{1/4} \alpha)$ for testing $C_4$-freeness}\label{subsec:C4-ub}
\label{sec:C4}
% Let $L$ denote the subset of vertices with degree at most $\alpha/\eps$
% \commentt{DR: We have a basis for this in the old main file. Question: do we use random walks both when the two high-degree vertices have degree greater than $n^{1/2}$ and when at least one of them has lower degree?}




\ICALPCR{In this subsection we prove the more general (arboricity-dependent) form of the upper bound for testing $C_4$-freeness which is stated as Theorem~\ref{thm:C4-ub} in the introduction.}


Recall that the assumption on $\alpha$ is that it is an upper bound on the arboricity $arb(G)$ - hence, while it is known that for graphs with $arb(G)>\sqrt{n}$ there exists a $C_4$,  we cannot simply reject if we get $\alpha>n^{1/2}$ since it might be that $arb(G)<\sqrt{n}$ (and we want one-sided error). 
The above expression shows that in the  case that $\alpha>n^{1/2}$, the $n^{1/4}\alpha$ term is replaced by $n^{3/4}$ (and the additive $\alpha$ term is due to the edge sampling).

%\commentt{DR: the issue of knowing $\alpha$}
% Now in preliminaries

The algorithm referred to in Theorem~\ref{thm:C4-ub} is described next. 
%The algorithm and its analysis refer to two degree thresholds: 
%\commentt{check/set constants}

\medskip\noindent
\textsf{Test-$C_4$-freeness}($n,\epsilon,\alpha$)\label{alg:C4}
\begin{enumerate}
\item Let $\theta_0 = 4\alpha/\epsilon$,
% c_0\cdot \alpha/\epsilon$,
$\theta_1 = c_1\cdot \sqrt{n}/\eps$ (where %$c_0$ and 
$c_1$ will be determined subsequently) and $\theta_{min}=\min\{\theta_0, \theta_1\}$  (it is useful to read the algorithm while having in mind that $\theta_0\leq \theta_1$ (i.e., $\alpha = O(\sqrt{n})$) so that $\theta_{min}=\theta_0$).
%\commentt{Set $c_0$ to 4. If for some reason want to return, then note that there are places where we write $c_1/4$}
\label{step:set_params_C4}
\item Repeat the following $t=\Theta(1/\epsilon)$ times:
  \begin{enumerate}
     \item Select an edge $e$ by calling the procedure \SelectEdge($\alpha,\epsilon$), which is provided in Appendix~\ref{append:sample-edge}.  If it
     % fails to return an edge, then return `Accept' (and terminate).
     does not return an edge, then continue to the next iteration.
     \item Select an endpoint $v$ of $e$ by flipping a fair coin.
     \item \label{step:random-neighbors} If $d(v) \leq \theta_1$, then select $s_1=\Theta( \sqrt{d(v)/\eps})$ ~($=O(n^{1/4}/\eps)$) random neighbors of $v$, and for each neighbor $u$ such that  
     $d(u) \leq \theta_{min}$
     query all the neighbors of $u$. 
     \item \label{step:random-walk} Otherwise ($d(v) > \theta_1$), perform $s_2=\Theta(\sqrt{  (n\alpha/\theta_1)\log n}/\eps^2)$ 
      \;$(=\tilde{O}(n^{1/4}\alpha^{1/2}/\eps^2))$ 
     random walks of length $2$ starting from $v$. 
     \item If a $C_4$ is detected, then return it, `Reject' and terminate.
  \end{enumerate}
  \item Return `Accept'.
\end{enumerate}
We note that the algorithm can be unified/simplified so that it only performs random walks of length-2 
where the number of walks is $\Theta(n^{1/4}\alpha/\eps^2)$, %(i.e., Step~\ref{step:random-walk} is performed for every selected $v$ and not only if $d(v)>\theta_1$),
but then the analysis becomes slightly more complicated. %\commentt{I guess it should be with them $\min\{n^{3/4},n^{1/4}\alpha\}$}


We start by stating a claim concerning the procedure \SelectEdge\ -- where both the procedure and the proof of the claim are deferred to Appendix~\ref{append:sample-edge}. We then state and prove two additional claims that will be used in the proof of Theorem~\ref{thm:C4-ub}.

\begin{claim}\label{clm:select-edge} 
With probability at least $2/3$ the procedure \SelectEdge\ returns an edge. Conditioned on it returning an edge, each edge incident to a vertex with degree at most $\theta_0$ is returned with probability at least $1/(2m')$ and at most $1/m'$, where $m'$ is the number of edges incident to vertices with degree at most $\theta_0$.
\end{claim}


\iffalse
\commentt{DR: The first claim stated next holds as is in multigraphs (since we are relying on the existence of sufficiently many edge-disjoint $C_4$s). Should decide if to add this to the statement of the claim or note this when we want to use it for $C_8$. The second claim does not hold as is for multigraphs (where two length-2 walks that are different and end at the same vertex might not induce a $C_4$ (e.g., there are two edges between $v$ and $u$ and two edges between $u$ and $v'$). I thought that think we can slightly modify it to fix this by looking at a union of two bipartite graphs: $(L_1,R)$ and $(L_2,R)$, but no (9/4) I think this might not work out, at least not easily. }
\fi
%%%%%%%%%%%%%%

\begin{claim}\label{clm:C4-ub1}
Let $v$ be a vertex and let $C(v,\theta_{min})$ be a set of edge-disjoint $C_4$'s 
containing $v$ such that the neighbors of $v$ on these $C_4$s all have degree at most $\theta_{min}$, where $\theta_{min}$ is as defined in the algorithm.\footnote{Actually, we do not rely on the setting of $\theta_{min}$, so this claim holds for any threshold value.} Suppose that $|C(v,\theta_{min})|\geq 1$ and let $\eps' = |C(v,\theta_{min})|/d(v)$.
If we select $s = 16\sqrt{d(v)/\eps'}$ random neighbors of $v$, and for each selected neighbor $u$ such that {$d(u)\leq \theta_{min}$} we query all the neighbors of $u$,
then the probability that we obtain a $C_4$ is at least $9/10$.
\end{claim}

\begin{proof}
Let $E'(v)$ denote the set of edges incident to $v$ that participate in the set $C(v,\theta_{min})$. By the premise of the claim, $|E'(v)|/d(v)= 2|C(v,\theta_{min})|/d(v)= 2\eps'$. Let $s'$ be the number of neighbors  of $v$ that are incident to edges in $E'(v)$ among the $s$ selected random neighbors of $v$.
It holds that $\Ex[s']=2\eps' \cdot s$, and  by the multiplicative Chernoff bound, $s' \geq \eps' \cdot s$
with probability at least $1-e^{-\eps'\cdot s/4}$. 
We first show that this probability is at least $19/20$, and then condition on this event.
By the setting of $s=16\sqrt{d(v)/\eps'}$, it holds that $  \eps'\cdot s = 16\sqrt{\eps'\cdot d(v)},$ and by the setting of  $\eps'=|C(v,\theta)|/d(v)$, we get $\eps'\cdot s \geq 16\sqrt{|C(v, \theta_{min})|}\geq 16$. 
Therefore, with probability at least $19/20$, $s'>\eps'\cdot s= 16\sqrt{\eps' \cdot d(v)}$.
We condition on this event and consider only those $s'$ 
selected neighbors of $v$ that are endpoints of $E'(v)$.

For each $4$-cycle $\rho\in C(v,\theta_{min})$, let $u_1(\rho)$ and $u_2(\rho)$ be the two neighbors of $v$ on this $C_4$ (so that they are endpoints of edges in $E'(v)$). Since the $C_4$s in $C(v,\theta_{min})$ are edge-disjoint, these vertices are distinct. 
Observe that the $s'$ selected neighbors of $v$   are uniformly distributed in $\bigcup_{\rho\in C(v,\theta_{min})}\{u_1(\rho),u_2(\rho)\}$, and that $s' \geq 16\cdot\sqrt{|C(v,\theta_{min})|}$. Hence, by the ``birthday paradox'',  with high constant probability, the sample of neighbors of $v$ contains two vertices, $u_1(\rho)$, and $u_2(\rho)$
for some $\rho \in C(v,\theta_{min})$. 
Conditioned on this event, once the (at most $\theta_{min}$) neighbors of
$u_1(\rho)$ and $u_2(\rho)$ are queried, the four-cycle $\rho$ is observed.
\iffalse
For every $i,j \in [s']$, $i<j$, we define a Bernoulli random variable $\chi_{i,j}$ that is 1 if and only if the $i$th and the $j$th random walks (among the $s'$ considered) induce one of the $C_4$s in $C(v,\theta_0)$. We next show that we can apply Corollary~\ref{cor:cheb1} (with $s$ in that corollary set to $s'$ -- see Appendix~\ref{app:prob}) to get an upper bound on the probability that $\sum_{i,j \in [s'], i<j} \chi_{i,j} = 0$.
Since the probability of this event is minimized when all the neighbors of $v$ on the $C_4$s in $C(v,\theta_0)$ have degree exactly $\theta_0$ (rather than at most $\theta_0$), we henceforth assume that this is the case.
% , which in turn is an upper bound on the probability that no $C_4$ is detected.

The first condition in the corollary is clearly satisfied by the definition of the random variables, as any two walks are independent.
Now consider the second condition.
 Let $e$ denote the first edge in the $i$th length-2 walk. Then $e$ belongs to a unique 4-cycle in $C(v,\theta_0)$, denote it $(v,v_1, v_2, v_3)$ and assume that $e=(v,v_1)$. Hence, for $\chi_{i,j}$ to be equal $1$, it must be the case that
 \begin{itemize}
     \item the second step of the $i$th walk (which starts at $v_1$)  hits $v_2$ -- this happens with probability $1/\theta_0$
     \item the first step of the $j$th walk  (which starts at $v$) hits $v_3$ -- this happens with probability $1/|E'(v)|$ (due to the conditioning that all the walks start at an edge that belongs to $C(v,\theta_0)$.)
     \item the second step of the $j$th walk (which starts at $v_3$) hits $v_2$ - this again happens with probability $1/\theta_0$  
 \end{itemize}
Therefore, for any pair $i,j\in [s']$, $i<j$ we have that
\begin{equation}
\mu = \Pr[\chi_{i,j}=1] = \frac{1}{\theta_0}\cdot \frac{1}{|E'(v)|}\cdot \frac{1}{\theta_0}
     = \frac{1}{2\eps'\cdot d(v)\cdot (\theta_0)^2}\;.
\end{equation}
Since $s' \geq \eps' \cdot s \geq  c\sqrt{\eps' \cdot d(v)}\cdot \theta_0$, the third condition in Corollary~\ref{cor:cheb1}
is satisfied (for $c_2$ in that claim set to $c'/\sqrt{2}$).
% For any pair of such random walks, the probability that the union of their edges is a $C_4$ is at least $\frac{1}{\theta_0}\cdot \frac{1}{|E'(v)|}\cdot \frac{1}{\theta_0} \geq \frac{1}{\eps'\cdot d(v)\cdot (\theta_0)^2}$.
% Since the number of pairs of random walks is ${s' \choose 2} \geq (c^2/2)\cdot \eps' \cdot d(v)\cdot (\theta_0)^2$, the claim follows by Chebishev's inequality. \commentt{DR: to fill in details}

It remains to verify that the second condition holds. For any four indices $i_1,j_1,i_2,j_2\in [s']$, $i_1 < j_1$, $i_2 < j_2$ such that exactly two of the four indices are the same, 
\begin{equation}
\Pr[\chi_{i_1,j_1}=\chi_{i_2,j_2} = 1] = \frac{1}{\theta_0}\cdot \left(\frac{1}{|E'(v)|}\cdot \frac{1}{\theta_0} \right)^2
 = \mu^2 \cdot \theta_0\;
\end{equation}
Since $\mu=1/(2\eps'\cdot d(v)\cdot \theta_0)$, 
and by its definition, $\eps' \geq 1/d(v)$, %\talya{for $d(v)>1/2\eps'$}, 
we have that $\mu^{1/2} < 1/\theta_0$, and hence the second condition in Corollary~\ref{cor:cheb1} holds as well (for $c_1 =1$).
\fi
% By setting $c' = 1/100$, the current claim follows.
\end{proof}

\begin{claim}\label{clm:C4-ub2}
Let $G$ be a graph over $n$ vertices and $m$ edges, and let $\theta_1, \eps', \eps''$ be parameters.
% \talya{added $\eps',\eps''$}
Suppose that $G$ contains a bipartite subgraph $G' = (L,R,E(G'))$ such that every vertex in $R$ has degree at least $\theta_1$ in $G$. 
%\talya{shouldn't it be in $G'$?}. % and $|E(G')| \geq \beta m$.
%\commentt{I think it has to be in $G$ - in $G'$ we only know that they have degree at least $\eps'$ times their degree in $v$}
Let $v$ be a vertex in $R$ such that $v$ has at least $\eps' \cdot d(v)$ neighbors in $L$ where each  of these neighbors, $u$, has at least $\eps'' \cdot \max\{d(u),\frac{m}{n}\}$ neighbors in $R$. %\talya{check max}
If $\theta_1\geq 2\sqrt{n/(\eps'\cdot \eps'')}$ and we take $s_2 \geq \frac{32}{\eps'\cdot \eps''}\cdot \sqrt{2\log n\cdot |R|}$ random walks of length 2 from $v$ for a sufficiently large constant $c'$, then with probability at least $9/10$ we shall detect a $C_4$ in $G$.
\end{claim}

\begin{proof}
For a pair of vertices $v$ and $v'\neq v$ in $R$, let $\ell_2(v,v')$ be the number of length-2 paths between $v$ and $v'$, and let $\ell_2(v,R) = \sum_{v'\in R} \ell_2(v,v')$.  Consider taking two random length-2 walks from $v$, and let $\mathcal{E}_1$ be the event that both of them end at vertices in $R$. Let $\mathcal{E}_2$ be the event that these two paths are distinct and end at the same vertex. Then for each single vertex $v'\in R$, conditioned on $\mathcal{E}_1$, the probability that the two walks end at $v'$ is exactly $\frac{\ell_2(v,v')}{\ell_2(v,R)}\cdot   \frac{\ell_2(v,v')-1}{\ell_2(v,R)}$.
Therefore,
\begin{equation} \label{eq:col}
\Pr[\mathcal{E}_2\,| \, \mathcal{E}_1] = \sum_{v'\in R} \frac{\ell_2(v,v')}{\ell_2(v,R)}\cdot   \frac{\ell_2(v,v')-1}{\ell_2(v,R)}
=  \frac{1}{(\ell_2(v,R))^2} \cdot \sum_{v'\in R} (\ell_2(v,v'))^2 - \frac{1}{\ell_2(v,R)}\;.
\end{equation}
We would like to lower bound the above probability.
For the first term on the right-hand-side, by applying the Cauchy-Schwartz inequality we get that
\begin{equation}\label{eq:col1}
\frac{1}{(\ell_2(v,R))^2} \cdot \sum_{v'\in R} (\ell_2(v,v'))^2 \geq  \frac{1}{(\ell_2(v,R))^2} \cdot |R| \cdot \left( \frac{\ell_2(v,R))}{|R|}\right)^2 = \frac{1}{|R|}\;.
\end{equation}
By combining Equations~\eqref{eq:col} and~\eqref{eq:col1} we get that
$\Pr[\mathcal{E}_2\,| \, \mathcal{E}_1] \geq \frac{1}{|R|} - \frac{1}{\ell_2(v,R)}\;.
$
Since each vertex in $R$ has degree at least $\theta_1$, we have that $|R| \leq \frac{2m}{\theta_1}$. 
By the premise of the claim regarding $v$ and its neighbors, $v$ has $\eps' d(v)\geq \eps'\cdot \theta_1$ neighbors in $L$, and each of them has 
at least $\eps''\cdot (m/n)$ neighbors in $R$. Therefore,
\begin{equation}
\ell_2(v,R) \geq \eps'\cdot \theta_1\cdot \eps''\cdot \frac{m}{n} 
 \geq \frac{\eps'\cdot \eps''\cdot \theta_1^2 \cdot |R|}{2n}\geq 2|R|,
\end{equation}
where the last inequality is by the premise $\theta_1 \geq 2\sqrt{n/(\eps'\cdot \eps'')}$. Therefore,
%\begin{equation}\label{eq:sol-R}
$\Pr[\mathcal{E}_2\,| \, \mathcal{E}_1] \geq \frac{1}{2|R|} \;.$
%\end{equation}
So far we have shown that when taking two distinct random walks from $v$, and conditioned on them both ending at $R$ (the event $\mathcal{E}_1$), the two paths collide on the end vertex (and hence result in a $C_4$) with probability at least $1/2|R|$. We shall now prove, 
that when taking $s$ length-2 random walks from $v$, sufficiently many of them indeed end at $R$, and that with high probability, at least two of them collide, resulting in a $C_4$.

Consider first the event $\mathcal{E}_1$.
By the premise of the claim, $v$ has at least $\eps'\cdot d(v)$ neighbors in $L$, and each $u$ of them has at least $\eps'' \max\{d(u), m/n\}\geq \eps'' d(u)$ neighbors in $R$. Therefore, 
the probability that a single random walk from $v$ ends at $R$ is at least $\eps'\cdot \eps''$. 
Hence, if we take $s\geq  \frac{32}{\eps' \cdot \eps''}\sqrt{2\log n \cdot |R|}$ length-2 random walks from $v$, and let $s'$ denote the number of  walks that end at a vertex in $R$,  we have that $\Ex[s']=32\cdot \sqrt{2\log n \cdot |R|}$, and that with probability at least $9/10$, we have $s' \geq 16\cdot \sqrt{2\log n \cdot |R|}$. We henceforth condition on this event. 

 Let $\chi'_{i,j}$ denote the event that the $i$th and $j$th random walks among the ones that end at $R$ collide on the ending vertex (and thus result in a $C_4$). By the above discussion, we have that for a specific pair $i\neq j$, $\Pr[\chi'_{i,j}=1]\geq 1/2|R|$. We now  lower bound the probability that at least one pair of random walks from the $s'$ that end in $R$ detects a $C_4$, i.e. lower bound $\sum_{i,j\in[s']}\chi_{i,j}$,  using Corollary~\ref{cor:cheb1}. For that end  we also need to upper bound the variance of the sum. 
 
 %Consider the sum $\sum_{v'\in R} \frac{\ell_2(v,v')}{\ell_2(v,R)}\cdot   \frac{\ell_2(v,v')-1}{\ell_2(v,R)}$ from Equation~\eqref{eq:col1}, which we have shown is lower bounded by $\frac{1}{2|R|}$. W
 Partition the vertices in $R$ according to $\ell_2(v,v')$, where $R_x(v) = \{v': 2^{x-1} < \ell_2(v,v') \leq 2^x\}$
for $x = 0, \dots \log L \leq \log n$.
Since $\sum_{v'\in R} \frac{\ell_2(v,v')}{\ell_2(v,R)}\cdot   \frac{\ell_2(v,v')-1}{\ell_2(v,R)}>\frac{1}{2|R|}$,  there exists at least one setting of $x$ for which
$\sum_{v'\in R_x} \frac{\ell_2(v,v')}{\ell_2(v,R)}\cdot   \frac{\ell_2(v,v')-1}{\ell_2(v,R)} \geq \frac{1}{2|R|\log n}$. We denote this setting by $x^*$ and observe that $x^* > 0$ (since for every $v'\in R_0$, $\ell_2(v,v')-1 =0$).

For every $i,j \in [s']$, $i<j$, we define a Bernoulli random variable $\chi_{i,j}$ that is 1 if and only if the $i$th and the $j$th random walks from $v$ (among the $s'$ considered) end at the same $v'\in R_{x^*}$ and pass through a different vertex in $L$. We next show that we can apply Corollary~\ref{cor:cheb1} (with $s$ in that claim set to $s'$) to get an upper bound on the probability that $\sum_{i,j \in [s'], i<j} \chi_{i,j} = 0$ (which is an upper bound on the probability that we do not detect a $C_4$).
% Since the probability of this event is minimized when all the neighbors of $v$ on the $C_4$s in $C(v,\theta_0)$ have degree exactly $\theta_0$ (rather than at most $\theta_0$), we henceforth assume that this is the case.
% , which in turn is an upper bound on the probability that no $C_4$ is detected.

%As in the proof of Claim~\ref{clm:C4-ub1},
 By the definition of the random variables, for every $i_1\neq i_2, j_1\neq j_2$, it holds that $\chi_{i_1,j_1}, \chi_{i_2, j_2}$ are independent,  so that 
the first condition in Corollary~\ref{cor:cheb1} is satisfied. Next, for any pair $i,j\in [s']$, $i<j$ we have that
\begin{equation} \label{eq:mu_clm6}
\mu = \Pr[\chi_{i,j}=1] = \sum_{v'\in R_{x^*}} \frac{\ell_2(v,v')}{\ell_2(v,R)}\cdot   \frac{\ell_2(v,v')-1}{\ell_2(v,R)} \geq \frac{1}{2|R|\log n}\;.
\end{equation}
Therefore, we have that  $s' \geq 16\cdot\sqrt{2|R|\log n}=16/\sqrt{\mu}$, and so  
the third condition in Corollary~\ref{cor:cheb1}
is satisfied (for $c_2=16$, where $s'$ serves as the parameter $s$ in the corollary).  

It remains to verify that the second condition holds. For any four indices $i_1,j_1,i_2,j_2\in [s']$, $i_1 < j_1$, $i_2 < j_2$ such that exactly two of the four indices are the same,
\begin{equation}
\Pr[\chi_{i_1,j_1}=\chi_{i_2,j_2} = 1] = \sum_{v'\in R_{x^*}} \frac{\ell_2(v,v')}{\ell_2(v,R)}\cdot   \left(\frac{\ell_2(v,v')-1}{\ell_2(v,R)}\right)^2 \leq \mu\cdot \frac{2^{x^*}-1}{\ell_2(v,R)} \;.
\end{equation}
Since by Equation~\eqref{eq:mu_clm6}, $\mu = \sum_{v'\in R_{x^*}} \frac{\ell_2(v,v')}{\ell_2(v,R)}\cdot   \frac{\ell_2(v,v')-1}{\ell_2(v,R)} \geq \frac{2^{2(x^*-1)}}{2(\ell_2(v,R))^2}$ (as $\ell_2(v,v') \geq 2^{x^*-1}$ for every $v'\in R_{x^*}$  and $\ell_2(v,v') - 1 \geq \ell_2(v,v')/2$), 
we get that $\Pr[\chi_{i_1,j_1}=\chi_{i_2,j_2} = 1]<\sqrt 2 \cdot \mu^{3/2}$, and so
the second condition in Corollary~\ref{cor:cheb1} holds as well (for $c_1=\sqrt 2$).
Thus, the current claim follows. 
\end{proof}





\medskip
We are now ready to prove Theorem~\ref{thm:C4-ub}.


\begin{proofof}{Theorem~\ref{thm:C4-ub}}
Since the algorithm only rejects a graph $G$ if it detects a $C_4$, it will always accept graphs that are $C_4$-free.
Hence, we focus on  the case that $G$ is $\eps$-far from being $C_4$-free.

%\commentt{changed a bit below, which then meant that ALL CONSTANTS in the first case had to be further modified (somehow, for the second case it was already stated for $|\calC_2|\geq \eps m/16$}
% We start by setting $c_0 = 4$ so that 
% $\theta_0 = 4\alpha/\epsilon$ (we later determine $c_1$, where recall that 
% $\theta_1 = c_1\sqrt{n}/\eps$).
Recall that $\theta_0 = 4\alpha/\epsilon$
and let $E_{>\theta_0}$ be the subset of edges in $G$ where both endpoints have degree greater than $\theta_0$. 
Since the arboricity of $G$ is at most $\alpha$,
there are at most $2m/\theta_0$ vertices with degree greater than $\theta_0$, so that $|E_{>\theta_0}| \leq (2m/\theta_0)\cdot \alpha=\eps m/2$ edges.

Since $G$ is $\eps$-far from $C_4$-free, if we remove all edges in $E_{>\theta_0}$, then we get a graph that is at least $(\eps/2)$-far from $C_4$-free. It follows that there exists a set of edge-disjoint $C_4$s, denoted $\calC$, such that
no $C_4$ in $\calC$ contains an edge in $E_{>\theta_0}$, and $|\calC| \geq \eps m/8$.

 We next  partition $\calC$ into two disjoint subsets: $\calC_1$ contains those $C_4$s that have at most one vertex with degree at least $\theta_1$ in them, and $\calC_2$ contains those that have at least two such vertices (where in the case $\theta_0\leq \theta_1$ there will be exactly two).  Since $\calC = \calC_1 \cup \calC_2$, either $|\calC_1| \geq \eps m/16$ or $|\calC_2| \geq \eps m/16$ (possibly both).
 % Since $\calC = \calC_{>\theta_0} \bigcup \calC_1 \bigcup \calC_2$, $|\calC|>\eps m$ and $|\calC_{>\theta_0}|<\eps m/2$, either $|\calC_1| \geq \eps m/4$ or $|\calC_2| \geq \eps m/4$. 

\paragraph{The case $\mathbf{|\calC_1| \geq \eps m/16}$.}
Consider first the case that $|\calC_1| \geq \eps m/16$. In order to analyze this case, we apply a process of ``coloring'' vertices and edges. Initially, all vertices and edges that participate in $C_4$s that belong to $\calC_1$ are colored \emph{green}, and all other vertices and edges are colored \emph{red}. We next apply the following iterative process. As long as there is a green vertex $v$ whose number of incident green edges is less than $\eps d(v)/64$, color $v$ and its green incident edges by red. Observe that the total number of edges colored red by this process is at most $\eps m/32$.
Furthermore, at the end of this process, every green vertex $v$ has at least $\eps d(v)/64$ incident green edges (and if a vertex is red, then all its incident edges are red).
Let $\calC_1'$ be the subset of $\calC_1$ that consists of those $C_4$s in $\calC_1$ whose edges all remain green after the process (and hence they are green), so that $|\calC'_1| \geq \eps m/32$.

By the definition of $\calC_1$, and hence also $\calC'_1$, in each $C_4$ in $\calC'_1$ there is at most  one vertex with degree greater than $\theta_1$, and no edges such that both endpoints have degree greater than $\theta_0$. Assume without loss of generality that for each four-cycle $\rho\in \calC'_1$, 
where $\rho=(v_0(\rho), v_1(\rho), v_2(\rho), v_3(\rho))$,  $v_2(\rho)$ is the highest degree vertex (where $d(v_2(\rho))$ could be any value between $1$ to $n$).  
Let $V_0(\calC'_1) =\bigcup_{\rho \in \calC'_1} \{v_0(\rho)\}$ denote this set of vertices (i.e., the ones that are across from the highest degree vertex in a (green) four-cycle in $\calC_1'$).
% \commentt{put in observation}

\medskip\noindent\textsf{Observation: }
% \begin{claim} \label{clm:deg_v0}
\emph{For every $\rho\in\calC'_1$,
\begin{inparaenum}[(1)]
    \item $d(v_0(\rho)) \leq \theta_1$, and
    \item  $v_1(\rho)$ and $v_3(\rho)$ are of degree at most $\theta_{min}=\min\{\theta_0, \theta_1\}$.
\end{inparaenum}
}
% \end{claim}

\smallskip
To verify this observation, note that
    by the definition of $\calC'_1$, for every $\rho\in \calC'_1$, there is at most one vertex with degree greater than $\theta_1$, and since $v_2(\rho)$ is the highest degree vertex in $\rho$, it follows  that all three other vertices in $\rho$ are of degree at most $\theta_1$.
    
    We now show that $d(v_1(\rho))\leq \theta_0$, and the proof for $v_3(\rho)$ is identical. If $d(v_2(\rho))>\theta_0$, then it 
    must be the case that $d(v_1(\rho))<\theta_0$, as otherwise both have degree greater than $\theta_0$  and so they cannot be connected, which is a contradiction to them both being incident on the four-cycle $\rho$.  
    %
    If  $d(v_2(\rho))\leq \theta_0$, then since  $v_2(\rho)$ is the highest degree vertex in $\rho$, $d(v_1(\rho))\leq d(v_2(\rho))\leq \theta_0$.

\medskip
Therefore, for every   $v\in V_0(\calC'_1)$, it has at least $\eps d(v)/64$ neighbors $u$ such that $(v,u)$ is green and  $d(u) \leq \theta_{min}$. Hence, overall in the graph,   the set of vertices $V_0(\calC'_1)$ has  at least $\eps m/32$ green edges that are incident to it and their second endpoint is  of degree at most $\theta_{min}\leq \theta_0$. It follows that     conditioned on an edge being returned by  procedure \SelectEdge, by Claim~\ref{clm:select-edge}, it returns an edge incident to a vertex $v\in V_0(\calC'_1)$ with probability at least $(\eps m/32)/2m'>\eps/128$ (since $m'>\frac{1}{2} m$).
So the probability that in some iteration of \TestCfour\ a vertex $v_0\in V_0(\calC'_1)$ is selected, is at least $1-(1-\frac{\eps}{128})^t >9/10$ (recall that $t = \Theta(1/\eps)$ so that it suffices to set $t= 500/\eps$). 


 Conditioning on this event, we apply Claim~\ref{clm:C4-ub1}. 
 Specifically:
 \begin{itemize}
     \item $\theta_0=4\alpha/\eps$ 
      (as defined in Step~\ref{step:set_params_C4} in Algorithm~\TestCfour); %with $c_0=4$);
     \item $C(v_0,\theta_{min})$ is the set of $C_4$s in $\calC'_1$ that are incident to $v_0$;
     \item $\eps' = |C(v_0,\theta_{min})|/d(v) \geq \eps/128$ (since $v_0$ has at least $\eps d(v)/64$ incident green edges, and they can be partitioned into pairs such that each pair belongs to exactly one $C_4$ in $C(v_0,\theta_{min})$);
     \item $d(v_0) \leq \theta_1$ (by 
     the above observation);
     %Claim~\ref{clm:deg_v0}).
 \end{itemize}
 %\end{inparaenum} 
 In order to apply the claim, we must ensure that $s>16\sqrt{d(v_0)/\eps'}$. By the above, it is sufficient to set $s_1$ in Step~\ref{step:random-neighbors}, to be $s_1=512\sqrt{d(v_0)/\eps}$.
 
 Hence, by Claim~\ref{clm:C4-ub1}, if Step~\ref{step:random-neighbors} is applied to $v_0$, then a $C_4$ is observed with probability at least $9/10$.

The analysis for the case that $\mathbf{|\calC_2| \geq \eps m/16}$ is similar, and due to space constraints, it is deferred to Appendix~\ref{append:case-C2}. 

We next turn to analyze the query complexity. By the settings of $\theta_0$, $\theta_1$, $t$, $s_1$ and $s_2$ in the algorithm,
% the two potential values of $s$ (either $64\sqrt{d(v)/\eps}=O(\theta_1/\eps)$ or $(10^5/\eps^2)\cdot \sqrt{2\log n (2n\alpha/\theta_1)}$), it holds that 
the query complexity of the algorithm is upper bounded as follows.
\begin{equation}
O\left(\frac{1}{\eps}\cdot \left(\frac{\alpha}{\eps}+\max\{s_1, s_2\}\right) \right) =
%
O\left(\frac{1}{\eps}\left( \frac{\alpha}{\eps}+\max\left\{\sqrt{\frac{\theta_1}{\eps}}\cdot \theta_{min}, \frac{1}{\eps^2}\cdot \sqrt{\frac{n\alpha}{\theta_1}\cdot\log n}\right\} \right) \right) 
\end{equation}
For the case that $\alpha\leq(c_1/4)\sqrt n$, we have that $\theta_{min}=\theta_0=\Theta(\alpha/\eps)$ and that $\theta_1=\Theta(\sqrt n/\eps)$, and so we get a complexity of
\begin{equation}
     O\left(
\eps^{-3} \cdot n^{1/4}\alpha^{1/2} \cdot \max\{\alpha^{1/2},\log^{1/2} n\}
 \right) = O(\eps^{-3}\cdot n^{1/4}\alpha)\;,
\end{equation}
where the last inequality is for $\alpha>\log n$, and otherwise the complexity is $O(\eps^{-3}\cdot n^{1/4}\alpha^{1/2}\log^{1/2} n)$. 

For the case that $\alpha>(c_1/4)\sqrt{n}$, we have that $\theta_{min}=\theta_1=\Theta(\sqrt n/\eps)$. Therefore, the complexity is
\begin{equation}
     O(\eps^{-3}\cdot (\alpha+n^{3/4}))\;.
\end{equation}
Thus, the proof is complete.
\end{proofof}



% \commentt{Should move the following claim to the appendix and just apply it here and in other places. One annoying thing is that it for $C_6$ we need a slightly different version with $s^2$ random variables}

% DR: moved following to preliminaries
%%%%%%%%%%%%%%%%%%%%%%%%%%
\iffalse
\begin{claim}\label{clm:cheb1}
Let $s$ be an integer, and denote by $\Phi(s)$ the set of pairs $i,j\in [s]$ such that $i<j$.
Let $\{\chi_{i,j}\}_{(i,j)\in \Phi(s)}$ be Bernoulli random variables where
 $\Pr[\chi_{i,j} = 1] = \mu$ for every $(i,j)\in \Phi(s)$. Suppose that the following conditions hold
 for some $c_1>0$ and $c_2 > 4$.
\begin{enumerate}
\item For every $(i_1,j_1)\in \Phi(s)$ and $(i_2,j_2)\in \Phi(s)$ such that the four indices are distinct, $\chi_{i_1,j_2}$ and $\chi_{i_2,j_2}$ are independent.
\item For every $(i_1,j_1)\in \Phi(s)$ and $(i_2,j_2)\in \Phi(s)$  such that exactly two of the four indices are the same,
% $|\{i_1,j_1,i_2,j_2\}|=3$,
$\Pr[\chi_{i_1,j_1}=\chi_{i_2,j_2} = 1] \leq c_1\cdot \mu^{3/2}$.
\item % For every $(i,j)\in \Phi(s)$, $\Pr[\chi_{i,j} = 1] = \mu$, where $\mu \geq \frac{1}{c s^2}$.
   $s \geq c_2/\sqrt{\mu}$
\end{enumerate}
Then $\Pr\left[ \sum_{(i,j)\in \Phi(s)} \chi_{i,j} = 0\right] \leq \frac{1+c_1}{c_2}$.
\end{claim}

\begin{proof}
For each pair $(i,j)\in \Phi(s)$,
let $\bar{\chi}_{i,j} = \chi_{i,j} - \mu$, so that $\bar{\chi}_{i,j} \leq \chi_{i,j}$, and
$\Ex[\bar{\chi}_{i,j}] = 0$.
By Chebishev's inequality,
\begin{equation}\label{eq:cheb1}
\Pr\left[\sum_{(i,j)\in \Phi(s)}\chi_{i,j}=0\right] \leq \frac{\Var\left[\sum_{(i,j)\in \Phi(s)} \chi_{i,j}\right]}{\left(\Ex\left[\sum_{(i,j)\in \Phi(s)}\chi_{i,j}\right]\right)^2}
= \frac{\sum_{i_1,j_2,i_2,j_2} \Ex[\bar{\chi}_{i_1,j_1}\cdot \bar{\chi}_{i_2,j_2}]}{{s\choose 2}^2\cdot \mu^2} \;.
\end{equation}
% we detect a triangle when selecting $\theta_1$ random neighbors of $a$ and $\theta_1$ random neighbors of $b$.
We upper bound the right-hand-side of Equation~\eqref{eq:cheb1} by breaking
 the sum $\sum_{i_1,j_2,i_2,j_2} \Ex[\bar{\chi}_{i_1,j_1}\cdot \bar{\chi}_{i_2,j_2}]$ into three sub-sums.

\paragraph{The first} is over $(i_1,j_1), (i_2,j_2)\in \Phi(s)$ such that all four indices are distinct. By the first condition in the claim, for such quadruples of indices, $\chi_{i_1,j_1}$ is independent of $\chi_{i_2,j_2}$ and so
\begin{equation}\label{eq:cheb-case1}
\Ex[\bar{\chi}_{i_1,j_1}\cdot \bar{\chi}_{i_2,j_2}] =  \Ex[\bar{\chi}_{i_1,j_1}]\cdot \Ex[\bar{\chi}_{i_2,j_2}] = 0\;.
\end{equation}
% where the last iequality follows from the definition of $\bar{\chi}_{i,j}$.
Hence, the first sub-sum does not contribute anything to the right-hand-side of Equation~\eqref{eq:cheb1}.

\paragraph{The second} is over $i_1 = i_2$ and $j_1 = j_2$, so that
\begin{equation}\label{eq:cheb-case2}
\Ex[\bar{\chi}_{i_1,j_1}\cdot \bar{\chi}_{i_2,j_2}] = \Ex[(\bar{\chi}_{i_1,j_1})^2]
\leq \Ex[(\chi_{i_1,j_1})^2] = \Ex[\chi_{i_1,j_1}] = \mu \;,
\end{equation}
where the second inequality follows from the definition of $\bar{\chi}_{i,j}$.
Since there are ${s\choose 2}$ quadruples in this sub-sum, after dividing by the denominator of the right-hand-side of Equation~\eqref{eq:cheb}, we get
\begin{equation}
\frac{{s\choose 2}\mu}{{s\choose 2}^2\mu^2} = \frac{1}{{s\choose 2} \mu}\;.
\end{equation}
By % Equation~\eqref{eq:mu-lb} and the setting of $\theta_1$,
the third condition in the claim,
this sub-sum contributes at most $\frac{4}{c_2^2}\leq \frac{1}{c_2}$.

\paragraph{The third} is over $(i_1,j_1), (i_2,j_2) \in \Phi(s)$ such that exactly two of the four indices are the same. By the second condition in the claim, for such quadruples,
\begin{equation}\label{eq:cheb-case3}
\Ex[\bar{\chi}_{i_1,j_1}\cdot \bar{\chi}_{i_2,j_2}]  =
% \Ex[\bar{\chi}_{i_1,j_1}\cdot \bar{\chi}_{i_1,j_2}]
\leq \Ex[\chi_{i_1,j_1}\cdot \chi_{i_2,j_2}] = \Pr[\chi_{i_1,j_1}=\chi_{i_2,j_2}=1] \leq c_1\mu^{3/2}\;.
\end{equation}
Since there are  ${s\choose 3}$ quadruples in this sub-sum, after dividing by the denominator of the right-hand-side of Equation~\eqref{eq:cheb} and using the  third conditions in the claim, we get that the contribution of this sub-sum is upper bounded by
\begin{equation}
{s\choose 3} \cdot \frac{c_1\mu^{3/2}}{{s\choose 2}^2 \mu^2} \leq \frac{c_1}{s\cdot \mu^{1/2}}\leq \frac{c_1}{c_2}\;.
\end{equation}

The claim follows by summing up the contributions of the three sub-sums.
\end{proof}
\fi
%%%%%%%%%%%%%%%%%%%%%%%%%%%%%



